\labeledsubsubsection{How to improve security?}
To overall increase security and protect against dependency-injection in any form, one must first of all secure their own systems.
Only proper permissions, file access and access-control will lead to a safe system.

Next, it is important to inform users of \textit{plugin-systems} or \textit{mods} that not every \textit{plugin}/\textit{mod} has necessarily good intensions.
Creating a central place for all \textit{mods} or \textit{plugins} to be uploaded to also increases the security.
Such a place must require a developer to sign-up so that, in case a malicious artifact was found, this user can be banned, the artifact removed and users informed about it.
Thus, a platform also has to have a comment or ranking system, including the option to report malicious uploads.\\
Even better, if the application has an integrated way of installing a \textit{plugin} or \textit{mod} in either the application itself or a launcher.
A good example for this is \texttt{Visual Studio Code}, which features an integrated and controlled marketplace for all sorts of \textit{plugins}.
Another example would be \texttt{NetBeans}, \texttt{IntelliJ IDEA/Jetbrains-Applications} or \texttt{Google Chrome/Chromium}.
All feature such a central hub for uploading and installing \textit{plugins}.

Finally, to be truly safe from dependency-injections one can execute an application only in an isolated, maybe containerized, environment with minimal to non additional tools and permissions.

Additionally, it may be possible to scan \textit{plugins} or \textit{mods} for malicious or unwanted code.
However, this process would be quiet expensive and probably also only works if the source code is available.